\begin{exo}
	Durant l'année scolaire 2018-2019, Ismaël a été lycéen en terminale, spécialité SVT.

	Il a utilisé ses moyennes annuelles pour faire une estimation du résultat qu'il pourrait obtenir au bac. Pour cela, il a utilisé les coefficients utilisés pour le bac cette année-là.
	\begin{center}
		\def\arraystretch{1.3}
		\begin{tabular}{|l|c|c|c|c|} \hline
			\textbf{Matière} & Maths      & Ph.-Ch.        & SVT    & Franç. \\ \hline
			\textbf{Coeff.}  & $7$        & $6$            & $8$    & $4$    \\ \hline
			\textbf{Moyenne} & $11,1$     & $12,3$         & $13,8$ & $10,5$ \\ \hline\hline
			\textbf{Matière} & Philo      & LV1            & LV2    & EPS    \\ \hline
			\textbf{Coeff.}  & $3$        & $3$            & $2$    & $2$    \\ \hline
			\textbf{Moyenne} & $10,4$     & $12,7$         & $9,2$  & $9,9$  \\ \hline\hline
			\textbf{Matière} & Hist.-Géo. & \textbf{Total} &        &        \\ \hline
			\textbf{Coeff.}  & $3$        &                &        &        \\ \hline
			\textbf{Moyenne} & $11,6$     &                &        &        \\ \hline
		\end{tabular}
	\end{center}

	Estimer la note qu'il a obtenu au bac sachant qu'il a été aussi stable qu'il l'a été durant l'année. Arrondir au centième.
\end{exo}
